% ============================================================
\chapter{Cha\^ines de Markov \`a Temps Continu}
\label{ch:markov_continu}
% ============================================================

Le passage du temps discret au temps continu n'est pas une simple formalit\'e~: il transforme profond\'ement la nature des cha\^ines de Markov. En temps discret, les transitions se produisent \`a des instants r\'eguliers. En temps continu, le processus peut sauter d'un \'etat \`a un autre \`a tout instant, et le temps de s\'ejour dans chaque \'etat suit une loi exponentielle --- la seule loi continue sans m\'emoire. Ce lien entre la propri\'et\'e de Markov et la loi exponentielle est l'une des observations les plus \'el\'egantes de la th\'eorie. Les CMTC mod\'elisent les files d'attente, les r\'eactions chimiques stochastiques, les r\'eseaux de t\'el\'ecommunications et bien d'autres syst\`emes o\`u les \'ev\'enements surviennent de mani\`ere al\'eatoire dans le temps.

\section{D\'efinition et propri\'et\'e de Markov}

\begin{definition}[Cha\^ine de Markov \`a temps continu]
Soit $E$ un espace d'\'etats d\'enombrable. Un processus $(X_t)_{t \geq 0}$
\`a valeurs dans $E$, \`a trajectoires c\`adl\`ag, est une \emph{cha\^ine de Markov
\`a temps continu} (CMTC) si pour tous $s, t \geq 0$ et tout $j \in E$~:
\[
  \PP(X_{t+s} = j \mid \mathcal{F}_s) = \PP(X_{t+s} = j \mid X_s) \quad \text{p.s.}
\]
La cha\^ine est \emph{homog\`ene} si $\PP(X_{t+s} = j \mid X_s = i) = p_{ij}(t)$
ne d\'epend pas de $s$.
\end{definition}

\begin{definition}[Semi-groupe de transition]
La famille de matrices $(P(t))_{t \geq 0}$ avec $P(t) = (p_{ij}(t))_{i,j \in E}$
est appel\'ee \emph{semi-groupe de transition}. Elle v\'erifie~:
\begin{enumerate}
  \item $P(0) = I$ (matrice identit\'e),
  \item $P(s+t) = P(s) P(t)$ pour tous $s, t \geq 0$ (propri\'et\'e de semi-groupe),
  \item $p_{ij}(t) \geq 0$ et $\sum_j p_{ij}(t) = 1$ pour tout $i$.
\end{enumerate}
\end{definition}

\section{G\'en\'erateur infinit\'esimal}

\begin{definition}[G\'en\'erateur infinit\'esimal]
Le \emph{g\'en\'erateur infinit\'esimal} (ou \emph{matrice des taux}) de la CMTC est
la matrice $Q = (q_{ij})_{i,j \in E}$ d\'efinie par~:
\[
  Q = \lim_{t \to 0^+} \frac{P(t) - I}{t}, \quad \text{c'est-\`a-dire} \quad
  q_{ij} = \lim_{t \to 0^+} \frac{p_{ij}(t) - \delta_{ij}}{t}.
\]
\end{definition}

\begin{proposition}[Propri\'et\'es de $Q$]
La matrice $Q$ v\'erifie~:
\begin{enumerate}
  \item $q_{ij} \geq 0$ pour $i \neq j$,
  \item $q_{ii} \leq 0$ pour tout $i$,
  \item $\sum_{j \in E} q_{ij} = 0$ pour tout $i$, c'est-\`a-dire $q_{ii} = -\sum_{j \neq i} q_{ij}$.
\end{enumerate}
On note $q_i = -q_{ii} = \sum_{j \neq i} q_{ij}$ le \emph{taux de sortie} de l'\'etat $i$.
Si $q_i > 0$, le temps de s\'ejour en $i$ suit une loi exponentielle
$\mathcal{E}(q_i)$.
\end{proposition}

\begin{intuition}[title=\textbf{Construction par temps de s\'ejour}]
Une CMTC peut \^etre construite comme suit~:
\begin{enumerate}
  \item Partant de l'\'etat $i$, le processus y reste un temps exponentiel de
    param\`etre $q_i$.
  \item \`A la fin de ce temps, il saute en $j \neq i$ avec probabilit\'e
    $q_{ij}/q_i$.
  \item Le processus recommence depuis le nouvel \'etat.
\end{enumerate}
La cha\^ine discr\`ete sous-jacente (les \'etats visit\'es successivement) a pour
matrice de transition $\tilde{P}$ avec $\tilde{p}_{ij} = q_{ij}/q_i$ pour $j \neq i$.
\end{intuition}

\section{\'Equations de Kolmogorov}

\begin{theoreme}[\'Equation de Kolmogorov progressive (forward)]
\label{thm:kolmo_forward}
Sous des conditions de r\'egularit\'e (par exemple $\sup_i q_i < \infty$)~:
\[
  P'(t) = P(t) Q, \quad \text{c'est-\`a-dire} \quad
  \frac{d}{dt} p_{ij}(t) = \sum_{k \in E} p_{ik}(t) q_{kj}.
\]
\end{theoreme}

\begin{theoreme}[\'Equation de Kolmogorov r\'etrograde (backward)]
\label{thm:kolmo_backward}
Sous les m\^emes conditions~:
\[
  P'(t) = Q P(t), \quad \text{c'est-\`a-dire} \quad
  \frac{d}{dt} p_{ij}(t) = \sum_{k \in E} q_{ik} p_{kj}(t).
\]
\end{theoreme}

\begin{remarque}
Si $E$ est fini, la solution est $P(t) = e^{Qt}$, o\`u l'exponentielle matricielle
est d\'efinie par $e^{Qt} = \sum_{n=0}^{\infty} \frac{(Qt)^n}{n!}$.
\end{remarque}

\begin{formules}[title=\textbf{\'Equations de Kolmogorov}]
\begin{align*}
  &\text{R\'etrograde (backward)}~: & P'(t) &= QP(t), \quad P(0) = I, \\
  &\text{Progressive (forward)}~: & P'(t) &= P(t)Q, \quad P(0) = I.
\end{align*}
Solution (espace d'\'etats fini)~: $P(t) = e^{Qt}$.
\end{formules}

\section{Mesures stationnaires et convergence}

\begin{definition}[Distribution stationnaire]
Un vecteur de probabilit\'e $\pi = (\pi_i)_{i \in E}$ est une \emph{distribution
stationnaire} pour la CMTC si~:
\[
  \pi Q = 0, \quad \text{c'est-\`a-dire} \quad \sum_{i \in E} \pi_i q_{ij} = 0
  \quad \forall j \in E.
\]
\'Equivalemment, $\pi P(t) = \pi$ pour tout $t \geq 0$.
\end{definition}

\begin{theoreme}[Convergence \`a l'\'equilibre]
\label{thm:convergence_equilibre}
Si la CMTC est irr\'eductible et r\'ecurrente positive, il existe une unique
distribution stationnaire $\pi$, et pour tous $i, j \in E$~:
\[
  \lim_{t \to \infty} p_{ij}(t) = \pi_j.
\]
\end{theoreme}

\section{Processus de naissance et de mort}

\begin{definition}[Processus de naissance et de mort]
Un processus de naissance et de mort sur $E = \N$ est une CMTC de g\'en\'erateur~:
\[
  q_{i,i+1} = \lambda_i \quad (\text{taux de naissance}), \quad
  q_{i,i-1} = \mu_i \quad (\text{taux de mort}), \quad i \geq 0,
\]
avec $\mu_0 = 0$, $\lambda_i > 0$ pour $i \geq 0$, $\mu_i > 0$ pour $i \geq 1$,
et $q_{ij} = 0$ pour $\abs{i - j} > 1$.
\end{definition}

\begin{theoreme}[Distribution stationnaire]
\label{thm:bd_stationnaire}
Le processus de naissance et de mort admet une distribution stationnaire si et
seulement si
\[
  S = \sum_{n=1}^{\infty} \frac{\lambda_0 \lambda_1 \cdots \lambda_{n-1}}
  {\mu_1 \mu_2 \cdots \mu_n} < \infty.
\]
Dans ce cas, $\pi_0 = (1 + S)^{-1}$ et
$\pi_n = \pi_0 \frac{\lambda_0 \cdots \lambda_{n-1}}{\mu_1 \cdots \mu_n}$.
\end{theoreme}

\begin{proof}
Les \'equations de bilan d\'etaill\'e $\pi_i q_{i,i+1} = \pi_{i+1} q_{i+1,i}$
donnent $\pi_i \lambda_i = \pi_{i+1} \mu_{i+1}$, d'o\`u la formule par r\'ecurrence.
La normalisabilit\'e \'equivaut \`a $S < \infty$.
\end{proof}

\begin{exemple}[File d'attente $M/M/1$]
Arriv\'ees selon un processus de Poisson de taux $\lambda$, services exponentiels
de taux $\mu$~: $\lambda_i = \lambda$, $\mu_i = \mu$ pour tout $i$. La condition
de stabilit\'e est $\rho = \lambda/\mu < 1$, et la distribution stationnaire est
g\'eom\'etrique~: $\pi_n = (1 - \rho)\rho^n$.
\end{exemple}

\begin{exemple}[File $M/M/c$]
Avec $c$ serveurs~: $\lambda_i = \lambda$ pour tout $i$ et
$\mu_i = \min(i, c) \mu$. La condition de stabilit\'e est $\rho = \lambda/(c\mu) < 1$.
\end{exemple}

\begin{exemple}[File $M/M/\infty$]
Nombre infini de serveurs~: $\lambda_i = \lambda$, $\mu_i = i\mu$.
La distribution stationnaire est de Poisson~: $\pi_n = e^{-\lambda/\mu}
\frac{(\lambda/\mu)^n}{n!}$.
\end{exemple}

\section{Explosion et r\'egularit\'e}

\begin{definition}[Explosion]
Une CMTC \emph{explose} si le nombre de sauts en temps fini est infini, c'est-\`a-dire
si $T_{\infty} = \sum_{n=0}^{\infty} S_n < \infty$ p.s., o\`u $S_n$ est le $n$-i\`eme
temps de s\'ejour.
\end{definition}

\begin{theoreme}[Crit\`ere de non-explosion]
\label{thm:non_explosion}
Si $\sup_{i \in E} q_i < \infty$ (taux de sortie born\'es), alors la CMTC
n'explose pas. Plus g\'en\'eralement, la cha\^ine n'explose pas si et seulement si
la solution minimale des \'equations de Kolmogorov r\'etrogrades est stochastique
(c'est-\`a-dire $\sum_j p_{ij}(t) = 1$).
\end{theoreme}

\begin{exemple}[Processus de Yule]
Le processus de naissance pure avec $\lambda_n = n\lambda$ (chaque individu se
reproduit ind\'ependamment au taux $\lambda$). On a $X_t = e^{\lambda t} X_0$
en moyenne, et la cha\^ine n'explose pas car
$\sum_n \frac{1}{q_n} = \sum_n \frac{1}{n\lambda} = +\infty$.
\end{exemple}

\section{Propri\'et\'e de Markov forte}

\begin{theoreme}[Propri\'et\'e de Markov forte --- temps continu]
\label{thm:markov_forte_continu}
Soit $(X_t)_{t \geq 0}$ une CMTC r\'eguli\`ere (sans explosion) et $\tau$ un
temps d'arr\^et fini p.s. Alors, conditionnellement \`a $X_{\tau} = i$, le processus
$(X_{\tau + t})_{t \geq 0}$ est une CMTC de m\^eme g\'en\'erateur, partant de $i$,
ind\'ependante de $\mathcal{F}_{\tau}$.
\end{theoreme}

\section{R\'eversibilit\'e en temps continu}

\begin{definition}[R\'eversibilit\'e]
Une CMTC de g\'en\'erateur $Q$ est r\'eversible par rapport \`a une distribution $\pi$
si~:
\[
  \pi_i q_{ij} = \pi_j q_{ji}, \quad \forall i, j \in E.
\]
\end{definition}

\begin{theoreme}[Th\'eor\`eme de Burke]
\label{thm:burke}
Dans une file $M/M/1$ en r\'egime stationnaire ($X_0 \sim \pi$), le processus
de d\'epart est un processus de Poisson de taux $\lambda$.
\end{theoreme}

\section{Exercices}

\begin{exercice}
Soit une CMTC sur $E = \{1, 2, 3\}$ de g\'en\'erateur
\[
  Q = \begin{pmatrix} -3 & 2 & 1 \\ 1 & -2 & 1 \\ 2 & 1 & -3 \end{pmatrix}.
\]
Calculer la distribution stationnaire. V\'erifier si la cha\^ine est r\'eversible.
\end{exercice}

\begin{exercice}
Pour un processus de naissance et de mort avec $\lambda_n = \lambda$ et
$\mu_n = n\mu$ ($n \geq 1$), trouver la distribution stationnaire et montrer
que la cha\^ine n'explose pas.
\end{exercice}

\begin{exercice}
Montrer que pour une CMTC sur un espace d'\'etats fini, l'exponentielle
matricielle $P(t) = e^{Qt}$ v\'erifie les \'equations de Kolmogorov
(progressive et r\'etrograde).
\end{exercice}

\begin{exercice}
Consid\'erer une file $M/M/1$ avec $\lambda = 2$ et $\mu = 3$. Calculer la
probabilit\'e que le syst\`eme soit vide en r\'egime stationnaire, le nombre moyen
de clients, et le temps moyen de s\'ejour (formule de Little).
\end{exercice}

\begin{exercice}
Montrer qu'un processus de naissance pure avec $\lambda_n = n^2$ explose en
temps fini. \emph{Indication~: montrer que $\sum_{n} 1/\lambda_n < \infty$.}
\end{exercice}

\begin{exercice}
Pour la cha\^ine d'Ehrenfest en temps continu (chaque boule est d\'eplac\'ee
ind\'ependamment au taux $1$), \'ecrire le g\'en\'erateur $Q$ et trouver la
distribution stationnaire.
\end{exercice}
